\section{Prediction-Powered Fusion for the ATE}
\label{sec:ate}

\subsection{Estimator and exact variance}

For fixed \((\lambda,\omega)\), define
\begin{align}
\widehat\theta_r(\lambda,\omega)
&=
\mathbb P_{R,n}Z_\lambda
+
\omega\left\{
\mathbb P_{O,N}\{r(X)g(X)\}
-
\mathbb P_{R,n}g(X)
\right\}.
\label{eq:ate-estimator}
\end{align}
Here \(Z_\lambda\) is the fusion score in
Definition~\ref{def:fusion-score}, \(g\) is defined in
Equation~\eqref{eq:obs-effect-prediction}, and \(r\) is the transport ratio in
Assumption~\ref{ass:exact-transport}. At the coefficient origin,
\begin{align*}
\widehat\theta_r(0,0)
&=\mathbb P_{R,n}Z_0
=\widehat\theta_{\mathrm{AIPW},R},
\end{align*}
so the proposed family contains the RCT-only AIPW estimator from
Definition~\ref{def:rct-only-aipw} exactly.
The first term averages an RCT-valid pseudo-outcome. Under Assumption~\ref{ass:exact-transport}, the expression in braces is a mean-zero control variate; under Assumption~\ref{ass:common-marginal}, this is the special case \(r=1\). Thus, the RCT identifies the target, \(\lambda\) changes within-covariate score noise, and \(\omega\) uses the large OBS covariate sample to remove predictable between-covariate variation.

\begin{theorem}[Exact mean and variance under honest evaluation]
\label{thm:ate-variance}
\label{thm:ate-exact-r}
Under Assumptions~\ref{ass:trial}, \ref{ass:exact-transport}, and
\ref{ass:honest-separation}, suppose \(n,N>0\) and the displayed scores are
square-integrable. Every fixed \((\lambda,\omega)\) then satisfies
\begin{align}
\mathbb E\{\widehat\theta_r(\lambda,\omega)\}
&=\theta,
\\
\operatorname{Var}\{\widehat\theta_r(\lambda,\omega)\}
&=
\frac1n\operatorname{Var}_R\{Z_\lambda-\omega g(X)\}
+
\frac{\omega^2}{N}\operatorname{Var}_O\{r(X)g(X)\}.
\label{eq:ate-variance}
\end{align}
Equivalently, the law of total variance gives the within-/between-covariate decomposition
\begin{align}
\operatorname{Var}\{\widehat\theta_r(\lambda,\omega)\}
&=
\frac1n\mathbb E_R\{\operatorname{Var}_R(Z_\lambda\mid X)\}
+\frac1n\operatorname{Var}_R\{\tau(X)-\omega g(X)\}
+\frac{\omega^2}{N}\operatorname{Var}_O\{r(X)g(X)\}.
\label{eq:ate-within-between}
\end{align}
\end{theorem}

The mean follows from \(\mathbb E_R(Z_\lambda)=\theta\) and \(\mathbb E_O(rg)=\mathbb E_R(g)\). The variance follows after rewriting Equation~\eqref{eq:ate-estimator} as the sum of independent RCT and OBS empirical means.

\begin{corollary}[Common-marginal specialization]
\label{cor:ate-common-active}
Under Assumptions~\ref{ass:trial}, \ref{ass:common-marginal}, and
\ref{ass:honest-separation}, set \(r=1\) in
Theorem~\ref{thm:ate-variance}. Then Equation~\eqref{eq:ate-estimator} is
exactly unbiased, and its OBS variance term is
\(\omega^2\operatorname{Var}_R\{g(X)\}/N\).
\end{corollary}

\subsection{Joint and population-orthogonal oracles}

Throughout this subsection, Assumption~\ref{ass:honest-separation} remains in
force, \(n\) and \(N\) are positive, and the displayed random variables have
finite second moments.

\begin{definition}[ATE variance-quadratic objects]
\label{def:ate-variance-objects}
Define the five population moments
\begin{align}
A&\triangleq\operatorname{Var}_R(\Delta),
&
B&\triangleq\operatorname{Var}_R(g)+\frac nN\operatorname{Var}_O(rg),
\\
C&\triangleq\operatorname{Cov}_R(Z_0,\Delta),
&
D&\triangleq\operatorname{Cov}_R(Z_0,g),
&
E&\triangleq\operatorname{Cov}_R(\Delta,g).
\label{eq:ate-moments}
\end{align}
The moment \(A\) is the curvature cost of changing the outcome-regression blend
\(\lambda\). The moment \(B\) is the curvature cost of the prediction-powered
scale \(\omega\), including both the RCT subtraction and the scaled OBS
addition. The moment \(C\) is the baseline-score covariance that determines the
linear slope in the \(\lambda\) direction. The moment \(D\) is the
baseline-prediction covariance that determines the linear slope in the
\(\omega\) direction, with a minus sign because the RCT score contains
\(-\omega g\). The moment \(E\) is the interaction covariance coupling the two
coefficients through the cross-term \(-2E\lambda\omega\).

Collect these moments and coefficients as
\begin{align}
\mathbf H_{\mathrm{ATE}}
&\triangleq
\begin{pmatrix}A&-E\\-E&B\end{pmatrix},
&
\boldsymbol\ell
&\triangleq
\begin{pmatrix}C\\-D\end{pmatrix},
&
\beta
&\triangleq
\begin{pmatrix}\lambda\\\omega\end{pmatrix},
\end{align}
\end{definition}

\begin{proposition}[ATE variance quadratic and positive-semidefinite curvature]
\label{prop:ate-variance-quadratic}
Under the standing conditions of this subsection, the compact variance in
Equation~\eqref{eq:ate-variance} is
\begin{align}
\operatorname{Var}\{\widehat\theta_r(\beta)\}
&=
\frac1n\left\{
\operatorname{Var}_R(Z_0)+2\boldsymbol\ell^\top\beta+\beta^\top\mathbf H_{\mathrm{ATE}}\beta
\right\}.
\label{eq:ate-quadratic}
\end{align}
Moreover, \(\mathbf H_{\mathrm{ATE}}\) is positive semidefinite because, for
every \((u,v)\in\mathbb R^2\),
\begin{align}
\begin{pmatrix}u&v\end{pmatrix}
\mathbf H_{\mathrm{ATE}}
\begin{pmatrix}u\\v\end{pmatrix}
&=
\operatorname{Var}_R(u\Delta-vg)
+\frac nN v^2\operatorname{Var}_O(rg)
\geq0.
\end{align}
\end{proposition}

The variance in Equation~\eqref{eq:ate-quadratic} contains the RCT-only
baseline \(\operatorname{Var}_R(Z_0)/n\), which no coefficient can change.
Coefficient selection can therefore focus on the remaining quadratic. Its
positive-semidefinite curvature guarantees that this reduced objective cannot
decrease without bound.

\begin{definition}[ATE coefficient objective]
\label{def:ate-coefficient-objective}
Define the coefficient-dependent part of the scaled variance by
\begin{align*}
\mathcal J_{\mathrm{ATE}}(\beta)
&\triangleq
2\boldsymbol\ell^\top\beta
+\beta^\top\mathbf H_{\mathrm{ATE}}\beta.
\end{align*}
\end{definition}

\begin{proposition}[Population orthogonality of the fusion directions]
\label{prop:ate-population-orthogonality}
Under Assumptions~\ref{ass:trial} and \ref{ass:honest-separation}, the two
endpoint scores have the same conditional target:
\begin{align*}
\mathbb E_R(Z_a\mid X)
&=\tau(X),
\qquad a\in\{0,1\},
\\
\mathbb E_R(\Delta\mid X)
&=0.
\end{align*}
Because \(g\) is a fixed function of \(X\), it follows that
\begin{align}
E
&=\operatorname{Cov}_R\{\Delta,g(X)\}=0,
&
D
&=\operatorname{Cov}_R\{\tau(X),g(X)\}.
\label{eq:ate-population-orthogonality}
\end{align}
\end{proposition}

This result formalizes the purpose of the outcome-regression blend: changing
\(\lambda\) changes pseudo-outcome noise without changing its conditional CATE
target. The zero population interaction separates the two oracle directions.
We nevertheless retain the general \(E\) term because its tuning-sample
analogue \(\widehat E\) is generally nonzero, and the empirical joint quadratic
program must account for that sampling fluctuation.

\begin{proposition}[Equivalent ATE oracle objectives]
\label{prop:ate-argmin-equivalence}
For every nonempty feasible set \(\mathcal B\subseteq\mathbb R^2\), the two
argmin sets, possibly empty when the infimum is not attained, satisfy
\begin{align*}
\operatorname*{arg\,min}_{\beta\in\mathcal B}
\operatorname{Var}\{\widehat\theta_r(\beta)\}
&=
\operatorname*{arg\,min}_{\beta\in\mathcal B}
\mathcal J_{\mathrm{ATE}}(\beta).
\end{align*}
\end{proposition}

Positive semidefiniteness means that the variance objective cannot curve
downward along any direction \((u,v)\). For this two-dimensional matrix, the
stronger condition \(AB-E^2>0\), together with positive semidefiniteness, makes
\(\mathbf H_{\mathrm{ATE}}\) positive definite and invertible. It turns a
possibly flat quadratic into one with a unique unrestricted minimizer.

\begin{definition}[Joint and ridge population oracles]
\label{def:ate-joint-ridge-oracles}
When \(AB-E^2>0\), define the unrestricted joint oracle by
\begin{align*}
\beta_{\mathrm{joint}}^\star
&\triangleq
\operatorname*{arg\,min}_{\beta\in\mathbb R^2}
\operatorname{Var}\{\widehat\theta_r(\beta)\}
=
\operatorname*{arg\,min}_{\beta\in\mathbb R^2}
\mathcal J_{\mathrm{ATE}}(\beta).
\end{align*}
The equality of the two argmin sets follows from
Proposition~\ref{prop:ate-argmin-equivalence}.
For a prespecified \(\rho>0\), define the ridge-stabilized population oracle by
\begin{align*}
\beta_\rho^\star
&\triangleq
\operatorname*{arg\,min}_{\beta\in\mathbb R^2}
\left\{\mathcal J_{\mathrm{ATE}}(\beta)+\rho\lVert\beta\rVert_2^2\right\}.
\end{align*}
Let \(I_2\) denote the \(2\times2\) identity matrix.
\end{definition}

\begin{theorem}[Joint and ridge population oracles]
\label{thm:joint-oracle}
If \(AB-E^2>0\), the joint oracle in
Definition~\ref{def:ate-joint-ridge-oracles} is unique and
\begin{align}
\beta_{\mathrm{joint}}^\star
&=
\begin{pmatrix}
\lambda_{\mathrm{joint}}^\star\\
\omega_{\mathrm{joint}}^\star
\end{pmatrix}
=-\mathbf H_{\mathrm{ATE}}^{-1}\boldsymbol\ell
=
\begin{pmatrix}
(ED-BC)/(AB-E^2)\\
(AD-EC)/(AB-E^2)
\end{pmatrix}.
\label{eq:joint-oracle}
\end{align}
For every \(\rho>0\), whether or not \(\mathbf H_{\mathrm{ATE}}\) is singular,
the ridge oracle is unique and
\begin{align}
\beta_\rho^\star
&=-(\mathbf H_{\mathrm{ATE}}+\rho I_2)^{-1}\boldsymbol\ell.
\label{eq:ridge-oracle}
\end{align}
\end{theorem}

Theorem~\ref{thm:joint-oracle} locates the joint minimum of the two-coefficient variance quadratic. Corollary~\ref{cor:joint-oracle-variance} answers the distinct comparison question: it evaluates the quadratic at that minimum and quantifies the improvement relative to the RCT-only origin \(\beta=0\). The general interaction \(E\) is retained here because its empirical analogue \(\widehat E\) need not vanish, even though the exact population model later implies \(E=0\).

\begin{remark}[Singular or constrained curvature]
If \(\mathbf H_{\mathrm{ATE}}\) is singular, an unpenalized joint minimum can
lie along a flat set rather than at a unique point. The Moore--Penrose and
null-space characterization, as well as existence and uniqueness results under
coefficient constraints, are recorded in Appendix~\ref{app:proofs}. The ridge
oracle remains unique for every \(\rho>0\). In implementation we minimize the
continuous quadratic over a prespecified compact box, which guarantees at least
one constrained solution.
\end{remark}

\begin{corollary}[Joint-oracle variance reduction]
\label{cor:joint-oracle-variance}
Let \(\sigma_0^2=\operatorname{Var}_R(Z_0)\) and \(\mathcal V_0=\sigma_0^2/n\) be the variance at \((\lambda,\omega)=(0,0)\). If \(AB-E^2>0\), the unrestricted joint oracle has
\begin{align*}
V_{\min}
&=
\mathcal V_0-
\frac1n
\frac{BC^2-2ECD+AD^2}{AB-E^2},
\\
\mathcal V_0-V_{\min}
&=
\frac1n\boldsymbol\ell^\top\mathbf H_{\mathrm{ATE}}^{-1}\boldsymbol\ell.
\end{align*}
The reduction is nonnegative and is strictly positive exactly when \(\boldsymbol\ell\neq0\), equivalently when \((C,D)\neq(0,0)\). Under population orthogonality \(E=0\) and \(A,B>0\), this specializes to
\begin{align*}
V_{\min}
&=
\mathcal V_0-
\frac1n\left(\frac{C^2}{A}+\frac{D^2}{B}\right).
\end{align*}
\end{corollary}

By Proposition~\ref{prop:ate-population-orthogonality}, the interaction term
\(-2E\lambda\omega\) vanishes. When \(A,B>0\), the population oracle therefore
separates:
\begin{align}
\lambda^\star
&=-\frac CA,
&
\omega^\star
&=\frac DB.
\label{eq:axis-oracle}
\end{align}
If \(A=0\), conditional mean zero forces \(\Delta=0\) almost surely and the \(\lambda\) direction is flat. If \(B=0\), both \(g(X)\) and \(r(X)g(X)\) are constant on their relevant supports, so the prediction-powered correction is samplewise zero and the \(\omega\) direction is flat. Under the common covariate marginal,
\begin{align*}
B
&=\left(1+\frac nN\right)\operatorname{Var}_R\{g(X)\}.
\end{align*}
For the rectangular box \(\mathcal B_\beta\) and positive \(A,B\), population orthogonality gives the coordinatewise projection
\begin{align*}
\lambda_{\mathcal B_\beta}^\star
&=\Pi_{[L_\lambda,U_\lambda]}(-C/A),
&
\omega_{\mathcal B_\beta}^\star
&=\Pi_{[L_\omega,U_\omega]}(D/B).
\end{align*}
For a coupled constraint or a nonzero interaction, the original quadratic must be minimized directly. The population oracle weakly improves the variance objective whenever \((0,0)\) is feasible. Estimated coefficients have no universal finite-sample no-harm guarantee.

\subsection{Estimable honest algorithm}

For random quantities \(U\) and \(W\) evaluated on
\(D_S^{\mathrm{tune}}\), assume \(n_S^{\mathrm{tune}}\geq2\). Conditional on
the fixed nuisance fits, let
\(\widehat{\operatorname{Cov}}_{S,\mathrm{tune}}(U,W)\) denote their usual
Bessel-corrected sample covariance and define
\(\widehat{\operatorname{Var}}_{S,\mathrm{tune}}(U)
\triangleq\widehat{\operatorname{Cov}}_{S,\mathrm{tune}}(U,U)\). The factor
\(n/N\) below uses the evaluation-sample sizes whose variance is optimized.

Using the tuning samples, compute
\begin{align}
\widehat A&=\widehat{\operatorname{Var}}_{R,\mathrm{tune}}(\Delta),
&
\widehat B&=\widehat{\operatorname{Var}}_{R,\mathrm{tune}}(g)
+\frac nN\widehat{\operatorname{Var}}_{O,\mathrm{tune}}(rg),
\\
\widehat C&=\widehat{\operatorname{Cov}}_{R,\mathrm{tune}}(Z_0,\Delta),
&
\widehat D&=\widehat{\operatorname{Cov}}_{R,\mathrm{tune}}(Z_0,g),
&
\widehat E&=\widehat{\operatorname{Cov}}_{R,\mathrm{tune}}(\Delta,g).
\end{align}
Define
\begin{align*}
\widehat{\mathbf H}_{\mathrm{ATE}}
&=
\begin{pmatrix}
\widehat A&-\widehat E\\
-\widehat E&\widehat B
\end{pmatrix},
&
\widehat{\boldsymbol\ell}
&=
\begin{pmatrix}\widehat C\\-\widehat D\end{pmatrix},
&
\widehat d
&=
\widehat A\widehat B-\widehat E^2.
\end{align*}
Choose positive thresholds \(t_A,t_B,t_{\det}\) before observing the tuning outcomes. The population-orthogonality axis rule is
\begin{align*}
\widehat\lambda_{\mathrm{axis}}
&=
\begin{cases}
-\widehat C/\widehat A,&\widehat A>t_A,\\
0,&\widehat A\leq t_A,
\end{cases}
&
\widehat\omega_{\mathrm{axis}}
&=
\begin{cases}
\widehat D/\widehat B,&\widehat B>t_B,\\
0,&\widehat B\leq t_B.
\end{cases}
\end{align*}
When \(\widehat A>t_A\), \(\widehat B>t_B\), and \(\widehat d>t_{\det}\), the unrestricted empirical joint rule is
\begin{align*}
\widehat\lambda_{\mathrm{joint}}
&=
\frac{\widehat E\widehat D-\widehat B\widehat C}{\widehat d},
&
\widehat\omega_{\mathrm{joint}}
&=
\frac{\widehat A\widehat D-\widehat E\widehat C}{\widehat d}.
\end{align*}
For the prespecified nonempty compact box \(\mathcal B_\beta=[L_\lambda,U_\lambda]\times[L_\omega,U_\omega]\) containing the origin, the boundary-correct empirical choice is
\begin{align*}
\widehat\beta_{\mathcal B_\beta}
&\in
\operatorname*{arg\,min}_{\beta\in\mathcal B_\beta}
\left\{
2\widehat{\boldsymbol\ell}^\top\beta
+\beta^\top\widehat{\mathbf H}_{\mathrm{ATE}}\beta
\right\}.
\end{align*}
If the unrestricted joint rule lies outside \(\mathcal B_\beta\), this quadratic program must be solved rather than clipping its coordinates separately; rectangular coordinate-wise clipping is exact only after imposing \(\widehat E=0\). For the joint rule, if a required threshold fails or the quadratic solver is numerically unstable, the prespecified fallback is \((0,0)\); the axis rule uses its coordinate-specific threshold definitions above.

The population-orthogonality plug-in imposes \(\widehat E=0\) and uses \((-\widehat C/\widehat A,\widehat D/\widehat B)\). The empirical joint minimizer retains \(\widehat E\) and applies Equation~\eqref{eq:joint-oracle}. The two rules differ because sample orthogonality is not exact.

Each sample covariance is conditionally unbiased for its corresponding
fixed-score population moment given the nuisance fits, but the ratios defining
\(\widehat\lambda\) and \(\widehat\omega\) are nonlinear and are not generally
unbiased for the population oracle coefficients. The thresholds and feasible
box control instability; they do not turn these selected coefficients into
finite-sample oracle estimates.

The displayed analytic ratios apply only when their corresponding denominator thresholds pass. Imposing \(\widehat E=0\) does not assert that the realized tuning covariance vanishes. It applies the proven population restriction \(E=0\) before solving the sample problem. Retaining \(\widehat E=\widehat{\operatorname{Cov}}_{R,\mathrm{tune}}(\Delta,g)\) instead minimizes the realized empirical quadratic, whose interaction term is generally nonzero because of tuning-sample noise.

Algorithm~\ref{alg:honest-ate} distinguishes the independent-role
implementation covered by the exact finite-sample results from optional nested
cross-fitting. Independent role samples may be supplied directly; physically
splitting one dataset is only one way to construct them.

\begin{algorithm}[H]
\caption{ATE estimation with honest evaluation and optional cross-fitting}
\label{alg:honest-ate}
\DontPrintSemicolon
\KwIn{Either role-indexed samples satisfying Assumption~\ref{ass:honest-separation} or prespecified paired outer folds; a feasible box and denominator thresholds.}
\textbf{Independent-role implementation:} Fit \(\widehat\mu_R\), \(\widehat\mu_O\), \(g\), and any ratio model on \(D_R^{\mathrm{nuis}},D_O^{\mathrm{nuis}}\). Compute the five moments and select the coefficients on \(D_R^{\mathrm{tune}},D_O^{\mathrm{tune}}\). Evaluate Equation~\eqref{eq:ate-estimator} on \(D_R^{\mathrm{eval}},D_O^{\mathrm{eval}}\). The exact theorem uses the true \(r\); replacing it by \(\widehat r\) invokes the drift conditions below.\;
\textbf{Optional cross-fitted implementation:} For each outer fold \(k\), fit nuisances and tune coefficients without that fold, score only its held-out RCT and OBS observations, and aggregate \(\widehat\theta_{\mathrm{CF}}=\sum_{k=1}^K a_k\widehat\theta_k\) using fixed \(a_k\geq0\) with \(\sum_k a_k=1\). Use \(a_k=1/K\) for equal-sized folds.\;
\textbf{Infer under independent roles:} Set \(\psi_R=Z_{\widehat\lambda}-\widehat\omega g\) and \(\psi_O=\widehat\omega rg\), and report
\(\widehat V=\widehat{\operatorname{Var}}_{R,\mathrm{eval}}(\psi_R)/n+\widehat{\operatorname{Var}}_{O,\mathrm{eval}}(\psi_O)/N\). Under conditional Lindeberg conditions, bounded conditional \((2+\delta)\) moments, and a nondegenerate variance limit, use
\(\widehat\theta_r\pm z_{1-\alpha/2}\sqrt{\widehat V}\).\;
\textbf{Cross-fitted inference caveat:} Do not sum foldwise variances or report a pooled analytic Wald interval without a separate stability or influence-function argument.\;
\KwOut{Independent roles: \(\widehat\theta_r\), \(\widehat V\), and its Wald interval; optional cross-fitting: \(\widehat\theta_{\mathrm{CF}}\) with the stated inference caveat.}
\end{algorithm}

For explicit evaluation notation, define
\begin{align*}
\psi_{R,i}
&=
Z_{\widehat\lambda}(V_{R,i}^{\mathrm{eval}})-\widehat\omega g(X_{R,i}^{\mathrm{eval}}),
&
\psi_{O,j}
&=
\widehat\omega r(X_{O,j}^{\mathrm{eval}})g(X_{O,j}^{\mathrm{eval}}),
\end{align*}
and let \(\widehat{\operatorname{Var}}_{R,\mathrm{eval}}\) and
\(\widehat{\operatorname{Var}}_{O,\mathrm{eval}}\) be the usual unbiased
sample variances with denominators \(n-1\) and \(N-1\), respectively. For
\(n,N\geq2\), the fixed-selected-coefficient variance estimator and Wald
interval are
\begin{align*}
\widehat V
&=
\frac{\widehat{\operatorname{Var}}_{R,\mathrm{eval}}(\psi_R)}{n}
+
\frac{\widehat{\operatorname{Var}}_{O,\mathrm{eval}}(\psi_O)}{N},
\\
\widehat\theta_r
&\mathbin{\pm}
z_{1-\alpha/2}\sqrt{\widehat V}.
\end{align*}
Unbiasedness of \(\widehat V\) follows from
Assumption~\ref{ass:honest-separation}. Wald coverage additionally requires the
stated Lindeberg, moment, and nondegeneracy conditions; it is not a finite-sample
coverage statement.
Here \(z_{1-\alpha/2}\) is the \((1-\alpha/2)\)-quantile of the standard normal distribution; for \(\alpha=0.05\), it is approximately \(1.96\).

\begin{proposition}[Exact variance estimation after honest selection]
\label{prop:ate-honest-variance-active}
Under Theorem~\ref{thm:ate-variance}, let \((\widehat\lambda,\widehat\omega)\)
be selected without the evaluation samples, as required by
Assumption~\ref{ass:honest-separation}, and suppose \(n,N\geq2\). The selected
estimator is unbiased and
\begin{align*}
\mathbb E(\widehat V)
&=
\operatorname{Var}\{\widehat\theta_r(\widehat\lambda,\widehat\omega)\}.
\end{align*}
\end{proposition}

\begin{theorem}[Final-sample central limit theorem and Wald inference]
\label{thm:ate-honest-clt-active}
In addition to Proposition~\ref{prop:ate-honest-variance-active}, suppose
\(n,N\to\infty\), \(n/N\to\kappa\in[0,\infty)\), and for some \(\delta>0\)
the centered RCT and OBS evaluation-score arrays have uniformly bounded
\((2+\delta)\) moments after the learned objects are fixed. This moment
assumption supplies the required Lindeberg condition. Assume also that
\begin{align*}
n\operatorname{Var}\{\widehat\theta_r(\widehat\lambda,\widehat\omega)\}
&\to\sigma^2\in(0,\infty)
\end{align*}
in probability. Then
\begin{align*}
\frac{\widehat\theta_r-\theta}{\sqrt{\widehat V}}
&\rightsquigarrow N(0,1),
\end{align*}
and the displayed Wald interval has asymptotic coverage \(1-\alpha\). Variance unbiasedness alone does not imply this coverage statement.
\end{theorem}

With an estimated ratio \(\widehat r\), define
\begin{align}
h_n(X)
&=\{\widehat r(X)-r(X)\}g(X).
\end{align}
The exact evaluation-sample drift is
\begin{align}
\omega b_{\widehat r}
&\triangleq
\omega\mathbb E_O\{h_n(X)\}.
\end{align}
Honest separation prevents evaluation leakage but does not remove this drift. A direct
sufficient negligible-error regime is
\begin{align}
\sqrt n\,\omega\,\mathbb E_O\{h_n(X)\}
&=o_p(1),
&
|\omega|\sqrt{\frac nN}
\left\|h_n-\mathbb E_O(h_n)\right\|_{L^2(P_O)}
&=o_p(1).
\end{align}
Evaluation-only variance is first-order adequate under these conditions and the
preceding CLT assumptions. Otherwise the ratio learner's influence-function
contribution, the OBS empirical-process term, and their covariances must be included.

\paragraph{Estimating the density ratio.}
All three constructions below use only the nuisance samples and are fixed
before coefficient tuning and final evaluation. Let
\(\mathbb P_R^{\mathrm{nuis}}\) and \(\mathbb P_O^{\mathrm{nuis}}\) denote the
corresponding empirical averages. First, pool the nuisance-sample covariates and
let \(S=1\) indicate an RCT
observation and \(S=0\) an OBS observation. With \(\pi=P(S=1)\) and
\(s(x)=P(S=1\mid X=x)\), Bayes' rule gives
\begin{align*}
r(x)
&=
\frac{1-\pi}{\pi}\frac{s(x)}{1-s(x)},
&
\widehat r_{\mathrm{cls}}(x)
&=
\frac{1-\widehat\pi}{\widehat\pi}
\frac{\widehat s(x)}{1-\widehat s(x)}.
\end{align*}
For an unweighted pooled classifier,
\(\widehat\pi=n_R^{\mathrm{nuis}}/(n_R^{\mathrm{nuis}}+n_O^{\mathrm{nuis}})\). Class weighting or resampling requires
replacing this factor by the effective training prior. The odds representation
also requires overlap, \(0<s(X)<1\), on the RCT covariate support. Probability
clipping stabilizes the estimate but changes it from the exact density ratio.

A Riesz construction estimates the transport functional directly \citep{chernozhukov2021riesz,chernozhukov2022autodml}. Let \(\phi:\mathcal X\to\mathbb R^p\) be a fixed feature map and define
\begin{align*}
\widehat M_O
&=
\mathbb P_O^{\mathrm{nuis}}\{\phi(X)\phi(X)^\top\},
&
\widehat b_R
&=
\mathbb P_R^{\mathrm{nuis}}\{\phi(X)\}.
\end{align*}
For \(\rho>0\), the empirical Riesz quadratic and its closed-form solution are
\begin{align*}
\widehat\eta_\rho
&=
\operatorname*{arg\,min}_{\eta\in\mathbb R^p}
\left[
\eta^\top\widehat M_O\eta
-2\widehat b_R^\top\eta
+\rho\lVert\eta\rVert_2^2
\right]
=
(\widehat M_O+\rho I_p)^{-1}\widehat b_R,
\\
\widehat r_{\mathrm{Riesz}}(x)
&=
\phi(x)^\top\widehat\eta_\rho.
\end{align*}
At \(\rho=0\), a compatible singular system uses \(\widehat M_O^+\widehat b_R\). A restricted feature map estimates a projected representer, and the resulting weights need not be nonnegative.

A nonnegative automatic-balancing alternative starts with an out-of-sample-evaluable family \(w_\eta:\mathcal X\to[0,\infty)\). For a fixed feature map \(\phi\), a positive semidefinite matrix \(W\), and \(\gamma>0\), define \(\lVert v\rVert_W^2=v^\top Wv\) and solve
\begin{align*}
\widehat\eta_{\mathrm{bal}}
&\in
\operatorname*{arg\,min}_{\eta:\,
\mathbb P_O^{\mathrm{nuis}}w_\eta=1}
\left[
\left\|
\mathbb P_O^{\mathrm{nuis}}\{w_\eta(X)\phi(X)\}
-
\mathbb P_R^{\mathrm{nuis}}\{\phi(X)\}
\right\|_W^2
+
\gamma\mathbb P_O^{\mathrm{nuis}}\{w_\eta(X)^2\}
\right],
\\
\widehat r_{\mathrm{bal}}(x)
&=
w_{\widehat\eta_{\mathrm{bal}}}(x).
\end{align*}
Including \(g\) among the coordinates of \(\phi\) directly targets the transported moment used by the ATE correction. The dispersion penalty controls unstable weights but generally prevents exact empirical balance.

These constructions do not automatically inherit the exact-r theorem. For any fixed evaluation-independent weight function \(w\), replacing \(r\) by \(w\) gives
\begin{align*}
\mathbb E\{\widehat\theta_w(\lambda,\omega)\}
&=
\theta
+
\omega\{P_O(wg)-P_R(g)\}
\\
&=
\theta
+
\omega P_O\{(w-r)g\}.
\end{align*}
Thus exact unbiasedness holds for \(w=r\), or more generally for a
fixed \(w\) satisfying the exact population balance condition
\(P_O(wg)=P_R(g)\). Classifier, restricted Riesz, and automatic-balancing
estimates are otherwise in the estimated-ratio regime governed by the drift and
rate conditions above. The Auto-DML citation motivates the Riesz learning
criterion; it does not by itself make the present score Neyman-orthogonal to
ratio error or supply the ratio learner's first-order variance contribution.

\paragraph{Nested outer cross-fitting.}
As an optional alternative to independent role samples, observations can be
reused without scoring any observation with a nuisance fit trained on its own
outcome. Prespecify paired outer folds in both sources, or randomize the folds
independently and hold the realized assignments fixed. Within each outer
complement, use inner splits to fit nuisances and select
\((\widehat\lambda_k,\widehat\omega_k)\); then score only the untouched paired
folds. With fixed nonnegative weights \(a_k\) summing to one, define
\begin{align*}
\widehat\theta_k
&=
\mathbb P_{R,k}Z_{\widehat\lambda_k}^{(-k)}
+\widehat\omega_k
\left\{
\mathbb P_{O,k}(r g^{(-k)})
-\mathbb P_{R,k}g^{(-k)}
\right\},
&
\widehat\theta_{\mathrm{CF}}
&=\sum_k a_k\widehat\theta_k.
\end{align*}

\begin{proposition}[What outer-fold honesty proves]
\label{prop:ate-outer-unbiased-active}
Under Assumption~\ref{ass:trial} and either the common covariate marginal or
exact \(r\), suppose the outer-fold assignment is prespecified and
data-independent, or independently randomized with its realized randomness
held fixed. Each \(\widehat\theta_k\) is then conditionally unbiased given its
own outer complement, and \(\widehat\theta_{\mathrm{CF}}\) is unconditionally
unbiased. Overlapping outer training complements make the fold estimators
dependent, so these facts do not justify adding foldwise variances or using a
naive pooled Wald interval.
\end{proposition}

A pooled cross-fitted interval requires an additional stability or influence-function argument, for example convergence of fold-specific RCT and OBS scores to common deterministic limits in their respective \(L_2\) norms, uniform \((2+\delta)\) moments, and a nondegenerate two-sample limit. The present theory does not assert a general cross-fitted variance theorem without such conditions.

\paragraph{Diagnostics.}
An implementation should report split assignments, propensity bounds, candidate boxes, all denominator and determinant thresholds, coefficient fallbacks, ratio overlap and clipping, and the selected RCT-only comparator. Near-zero \(\widehat A\), \(\widehat B\), or \(\widehat d\) must trigger the prespecified rule rather than an unstable analytic ratio.
